
\documentclass{article}
\usepackage{colortbl}
\usepackage{makecell}
\usepackage{multirow}
\usepackage{supertabular}

\begin{document}

\newcounter{utterance}

\centering \large Interaction Transcript for game `dond', experiment `coop\_it', episode 16 with gpt{-}5{-}2025{-}08{-}07{-}no{-}reasoning{-}t1.0.
\vspace{24pt}

{ \footnotesize  \setcounter{utterance}{1}
\setlength{\tabcolsep}{0pt}
\begin{supertabular}{c@{$\;$}|p{.15\linewidth}@{}p{.15\linewidth}p{.15\linewidth}p{.15\linewidth}p{.15\linewidth}p{.15\linewidth}}
    \# & $\;$A & \multicolumn{4}{c}{Game Master} & $\;\:$B\\
    \hline

    \theutterance \stepcounter{utterance}  
    & & \multicolumn{4}{p{0.6\linewidth}}{
        \cellcolor[rgb]{0.9,0.9,0.9}{
            \makecell[{{p{\linewidth}}}]{
                \texttt{\tiny{[P1$\langle$GM]}}
                \texttt{State giocando a un gioco di negoziazione cooperativa in cui dovete accordarvi su come dividere una serie di oggetti tra voi e un altro giocatore.} \\
\\ 
\texttt{Regole:} \\
\texttt{(a) A Lei e all'altro giocatore viene dato un insieme di oggetti. Ognuno di voi riceve anche una funzione di valore segreta, che rappresenta il valore di ciascun tipo di oggetto.} \\
\texttt{(b) Si scambiano messaggi con l'altro giocatore per concordare chi si aggiudica gli oggetti. Potete inviare un massimo di 5 messaggi ciascuno, oppure terminare in anticipo facendo una proposta segreta in qualsiasi momento.} \\
\texttt{(c) A ciascuno di voi viene chiesto di inviare una proposta segreta indicando gli oggetti che desiderate, formattata tra parentesi quadre come segue: "{[}Proposta: <numero> <nome oggetto>, <numero> <nome oggetto>, <...>{]}".} \\
\texttt{(d) Se le vostre proposte sono complementari, cioè ci sono abbastanza oggetti per soddisfare entrambe le proposte, a ciascun giocatore viene assegnato un punteggio basato sulla somma dei valori degli oggetti ricevuti. In caso contrario, entrambi ricevono zero punti.} \\
\texttt{(e) L'obiettivo è massimizzare la somma del proprio punteggio e di quello dell'altro giocatore.} \\
\\ 
\texttt{Cominciamo.} \\
\\ 
\texttt{L'insieme degli oggetti disponibili è:} \\
\\ 
\texttt{1 secchio, 5 libri, 1 specchio} \\
\\ 
\texttt{La funzione valore segreta è:} \\
\\ 
\texttt{secchio: 4, libro: 1, specchio: 1} \\
\\ 
\texttt{IMPORTANTE: i vostri messaggi, a meno che non si tratti di una proposta segreta, vengono trasmessi direttamente all'altro giocatore, quindi non includete alcuna risposta alle regole o testo di annuncio del vostro messaggio. Per fare una proposta segreta, utilizzate il formato indicato. Non utilizzare delle parentesi quadre quando si comunica all'altro giocatore, altrimenti verrà interpretata come una proposta segreta.} \\
            }
        }
    }
    & \\ \\

    \theutterance \stepcounter{utterance}  
    & \multicolumn{4}{p{0.6\linewidth}}{
        \cellcolor[rgb]{0.9,0.9,0.9}{
            \makecell[{{p{\linewidth}}}]{
                \texttt{\tiny{[P1$\rangle$GM]}}
                \texttt{Ciao! Grazie per giocare in modo cooperativo. Ecco i miei valori: secchio 4, ogni libro 1, specchio 1. A me interessa soprattutto il secchio; i libri e lo specchio hanno valore minore.} \\
\\ 
\texttt{Proposta cooperativa iniziale: tu prendi lo specchio e 2 libri; io prendo il secchio e 3 libri. Questo assegna tutto e dà un buon valore totale. Che ne pensi? Se hai valori diversi, dimmi quali oggetti per te valgono di più così possiamo ottimizzare la somma.} \\
            }
        }
    }
    & & \\ \\

    \theutterance \stepcounter{utterance}  
    & & & \multicolumn{2}{p{0.3\linewidth}}{
        \cellcolor[rgb]{0.9,0.9,0.9}{
            \makecell[{{p{\linewidth}}}]{
                \texttt{\tiny{[GM$|$GM]}}
                \texttt{continue} \\
            }
        }
    }
    & & \\ \\

    \theutterance \stepcounter{utterance}  
    & & & \multicolumn{2}{p{0.3\linewidth}}{
        \cellcolor[rgb]{0.9,0.9,0.9}{
            \makecell[{{p{\linewidth}}}]{
                \texttt{\tiny{[GM$|$GM]}}
                \texttt{Ciao! Grazie per giocare in modo cooperativo. Ecco i miei valori: secchio 4, ogni libro 1, specchio 1. A me interessa soprattutto il secchio; i libri e lo specchio hanno valore minore.} \\
\\ 
\texttt{Proposta cooperativa iniziale: tu prendi lo specchio e 2 libri; io prendo il secchio e 3 libri. Questo assegna tutto e dà un buon valore totale. Che ne pensi? Se hai valori diversi, dimmi quali oggetti per te valgono di più così possiamo ottimizzare la somma.} \\
            }
        }
    }
    & & \\ \\

    \theutterance \stepcounter{utterance}  
    & & \multicolumn{4}{p{0.6\linewidth}}{
        \cellcolor[rgb]{0.9,0.9,0.9}{
            \makecell[{{p{\linewidth}}}]{
                \texttt{\tiny{[GM$\rangle$P2]}}
                \texttt{State giocando a un gioco di negoziazione cooperativa in cui dovete accordarvi su come dividere una serie di oggetti tra voi e un altro giocatore.} \\
\\ 
\texttt{Regole:} \\
\texttt{(a) A Lei e all'altro giocatore viene dato un insieme di oggetti. Ognuno di voi riceve anche una funzione di valore segreta, che rappresenta il valore di ciascun tipo di oggetto.} \\
\texttt{(b) Si scambiano messaggi con l'altro giocatore per concordare chi si aggiudica gli oggetti. Potete inviare un massimo di 5 messaggi ciascuno, oppure terminare in anticipo facendo una proposta segreta in qualsiasi momento.} \\
\texttt{(c) A ciascuno di voi viene chiesto di inviare una proposta segreta indicando gli oggetti che desiderate, formattata tra parentesi quadre come segue: "{[}Proposta: <numero> <nome oggetto>, <numero> <nome oggetto>, <...>{]}".} \\
\texttt{(d) Se le vostre proposte sono complementari, cioè ci sono abbastanza oggetti per soddisfare entrambe le proposte, a ciascun giocatore viene assegnato un punteggio basato sulla somma dei valori degli oggetti ricevuti. In caso contrario, entrambi ricevono zero punti.} \\
\texttt{(e) L'obiettivo è massimizzare la somma del proprio punteggio e di quello dell'altro giocatore.} \\
\\ 
\texttt{Cominciamo.} \\
\\ 
\texttt{L'insieme degli oggetti disponibili è:} \\
\\ 
\texttt{1 secchio, 5 libri, 1 specchio} \\
\\ 
\texttt{La funzione valore segreta è:} \\
\\ 
\texttt{secchio: 2, libro: 0, specchio: 8} \\
\\ 
\texttt{IMPORTANTE: i vostri messaggi, a meno che non si tratti di una proposta segreta, vengono trasmessi direttamente all'altro giocatore, quindi non includete alcuna risposta alle regole o testo di annuncio del vostro messaggio. Per fare una proposta segreta, utilizzate il formato indicato. Non utilizzare delle parentesi quadre quando si comunica all'altro giocatore, altrimenti verrà interpretata come una proposta segreta.} \\
\\ 
\\ 
\texttt{Ciao! Grazie per giocare in modo cooperativo. Ecco i miei valori: secchio 4, ogni libro 1, specchio 1. A me interessa soprattutto il secchio; i libri e lo specchio hanno valore minore.} \\
\\ 
\texttt{Proposta cooperativa iniziale: tu prendi lo specchio e 2 libri; io prendo il secchio e 3 libri. Questo assegna tutto e dà un buon valore totale. Che ne pensi? Se hai valori diversi, dimmi quali oggetti per te valgono di più così possiamo ottimizzare la somma.} \\
            }
        }
    }
    & \\ \\

    \theutterance \stepcounter{utterance}  
    & & & \multicolumn{4}{p{0.6\linewidth}}{
        \cellcolor[rgb]{0.9,0.9,0.9}{
            \makecell[{{p{\linewidth}}}]{
                \texttt{\tiny{[GM$\langle$P2]}}
                \texttt{Grazie per la trasparenza! I miei valori sono: secchio 2, ogni libro 0, specchio 8. Quindi per me conta quasi solo lo specchio; i libri non valgono nulla e il secchio vale poco.} \\
\\ 
\texttt{La tua proposta iniziale (tu secchio+3 libri, io specchio+2 libri) per me va benissimo: io ottengo il massimo (specchio) e non mi interessa dei libri. La somma totale sarebbe: tu 4 + 3 = 7, io 8 = 8, totale 15.} \\
\\ 
\texttt{Se vuoi massimizzare ulteriormente il tuo punteggio senza ridurre il totale, posso lasciarti tutti i libri: io prendo solo lo specchio, tu secchio e 5 libri. Totale resta 12 (tu) + 8 (io) = 20? Correggo: con i tuoi valori, secchio 4 + 5 libri (5) = 9 per te; io 8 per lo specchio; totale 17, che è meglio di 15.} \\
\\ 
\texttt{Propongo quindi: tu prendi secchio e 5 libri; io prendo lo specchio. Se ti va, confermiamo. Se preferisci, posso anche accettare subito la tua prima proposta.} \\
            }
        }
    }
     \\ \\

    \theutterance \stepcounter{utterance}  
    & & & \multicolumn{2}{p{0.3\linewidth}}{
        \cellcolor[rgb]{0.9,0.9,0.9}{
            \makecell[{{p{\linewidth}}}]{
                \texttt{\tiny{[GM$|$GM]}}
                \texttt{continue} \\
            }
        }
    }
    & & \\ \\

    \theutterance \stepcounter{utterance}  
    & & & \multicolumn{2}{p{0.3\linewidth}}{
        \cellcolor[rgb]{0.9,0.9,0.9}{
            \makecell[{{p{\linewidth}}}]{
                \texttt{\tiny{[GM$|$GM]}}
                \texttt{Grazie per la trasparenza! I miei valori sono: secchio 2, ogni libro 0, specchio 8. Quindi per me conta quasi solo lo specchio; i libri non valgono nulla e il secchio vale poco.} \\
\\ 
\texttt{La tua proposta iniziale (tu secchio+3 libri, io specchio+2 libri) per me va benissimo: io ottengo il massimo (specchio) e non mi interessa dei libri. La somma totale sarebbe: tu 4 + 3 = 7, io 8 = 8, totale 15.} \\
\\ 
\texttt{Se vuoi massimizzare ulteriormente il tuo punteggio senza ridurre il totale, posso lasciarti tutti i libri: io prendo solo lo specchio, tu secchio e 5 libri. Totale resta 12 (tu) + 8 (io) = 20? Correggo: con i tuoi valori, secchio 4 + 5 libri (5) = 9 per te; io 8 per lo specchio; totale 17, che è meglio di 15.} \\
\\ 
\texttt{Propongo quindi: tu prendi secchio e 5 libri; io prendo lo specchio. Se ti va, confermiamo. Se preferisci, posso anche accettare subito la tua prima proposta.} \\
            }
        }
    }
    & & \\ \\

    \theutterance \stepcounter{utterance}  
    & & \multicolumn{4}{p{0.6\linewidth}}{
        \cellcolor[rgb]{0.9,0.9,0.9}{
            \makecell[{{p{\linewidth}}}]{
                \texttt{\tiny{[P1$\langle$GM]}}
                \texttt{Grazie per la trasparenza! I miei valori sono: secchio 2, ogni libro 0, specchio 8. Quindi per me conta quasi solo lo specchio; i libri non valgono nulla e il secchio vale poco.} \\
\\ 
\texttt{La tua proposta iniziale (tu secchio+3 libri, io specchio+2 libri) per me va benissimo: io ottengo il massimo (specchio) e non mi interessa dei libri. La somma totale sarebbe: tu 4 + 3 = 7, io 8 = 8, totale 15.} \\
\\ 
\texttt{Se vuoi massimizzare ulteriormente il tuo punteggio senza ridurre il totale, posso lasciarti tutti i libri: io prendo solo lo specchio, tu secchio e 5 libri. Totale resta 12 (tu) + 8 (io) = 20? Correggo: con i tuoi valori, secchio 4 + 5 libri (5) = 9 per te; io 8 per lo specchio; totale 17, che è meglio di 15.} \\
\\ 
\texttt{Propongo quindi: tu prendi secchio e 5 libri; io prendo lo specchio. Se ti va, confermiamo. Se preferisci, posso anche accettare subito la tua prima proposta.} \\
            }
        }
    }
    & \\ \\

    \theutterance \stepcounter{utterance}  
    & \multicolumn{4}{p{0.6\linewidth}}{
        \cellcolor[rgb]{0.9,0.9,0.9}{
            \makecell[{{p{\linewidth}}}]{
                \texttt{\tiny{[P1$\rangle$GM]}}
                \texttt{Perfetto, la tua proposta massimizza la somma totale: tu prendi secchio e 5 libri (valore per te 9), io prendo lo specchio (valore per me 8). Totale 17.} \\
\\ 
\texttt{Procedo con la proposta segreta complementare alla tua:} \\
\\ 
\texttt{{[}Proposta: 1 specchio{]}} \\
            }
        }
    }
    & & \\ \\

    \theutterance \stepcounter{utterance}  
    & & & \multicolumn{2}{p{0.3\linewidth}}{
        \cellcolor[rgb]{0.9,0.9,0.9}{
            \makecell[{{p{\linewidth}}}]{
                \texttt{\tiny{[GM$|$GM]}}
                \texttt{proposal} \\
            }
        }
    }
    & & \\ \\

    \theutterance \stepcounter{utterance}  
    & & & \multicolumn{2}{p{0.3\linewidth}}{
        \cellcolor[rgb]{0.9,0.9,0.9}{
            \makecell[{{p{\linewidth}}}]{
                \texttt{\tiny{[GM$|$GM]}}
                [0, 0, 1]
            }
        }
    }
    & & \\ \\

    \theutterance \stepcounter{utterance}  
    & & \multicolumn{4}{p{0.6\linewidth}}{
        \cellcolor[rgb]{0.9,0.9,0.9}{
            \makecell[{{p{\linewidth}}}]{
                \texttt{\tiny{[GM$\rangle$P2]}}
                \texttt{L'altro giocatore ha presentato una proposta segreta. Ora tocca a Lei presentare una sua proposta.} \\
            }
        }
    }
    & \\ \\

    \theutterance \stepcounter{utterance}  
    & & & \multicolumn{4}{p{0.6\linewidth}}{
        \cellcolor[rgb]{0.9,0.9,0.9}{
            \makecell[{{p{\linewidth}}}]{
                \texttt{\tiny{[GM$\langle$P2]}}
                \texttt{{[}Proposta: 1 secchio, 5 libri{]}} \\
            }
        }
    }
     \\ \\

    \theutterance \stepcounter{utterance}  
    & & & \multicolumn{2}{p{0.3\linewidth}}{
        \cellcolor[rgb]{0.9,0.9,0.9}{
            \makecell[{{p{\linewidth}}}]{
                \texttt{\tiny{[GM$|$GM]}}
                \texttt{proposal} \\
            }
        }
    }
    & & \\ \\

    \theutterance \stepcounter{utterance}  
    & & & \multicolumn{2}{p{0.3\linewidth}}{
        \cellcolor[rgb]{0.9,0.9,0.9}{
            \makecell[{{p{\linewidth}}}]{
                \texttt{\tiny{[GM$|$GM]}}
                [1, 5, 0]
            }
        }
    }
    & & \\ \\

    \theutterance \stepcounter{utterance}  
    & & & \multicolumn{2}{p{0.3\linewidth}}{
        \cellcolor[rgb]{0.9,0.9,0.9}{
            \makecell[{{p{\linewidth}}}]{
                \texttt{\tiny{[GM$|$GM]}}
                [[0, 0, 1], [1, 5, 0]]
            }
        }
    }
    & & \\ \\

\end{supertabular}
}

\end{document}
